\documentclass{tropical}
\title{Lifting the exponent}
\author{Nyan}
\date\today
\begin{document}
\begin{frame}{}
\maketitle
\end{frame}
\begin{frame}{Motivation}
(AIME 2020)\\
Let $n$ be the least positive integer for which 
$149^n - 2^n$ is divisible by $3^3 \cdot 5^5 \cdot 7^7$. 
Find the number of positive divisors of $n$.\\
Intuitively you just deal with the different prime powers
separately.
\end{frame}
\begin{frame}
\begin{thm}[Notation]
$\nu_p(n)$ denotes the exponent of $p$ in the prime factor of $n$.
For example, $\nu_2{1434}=1$.
\end{thm}
\begin{block}{LTE Statement}
Let $p$ be an odd prime and suppose $x,y$ are ints with $p\mid x-y$,
but not $p\mid x,y$. Then
\[\nu_p(x^n-y^n)=\nu_p(x-y)+\nu_p n\]
for all int $n$.
\end{block}
\end{frame}
\begin{frame}{Edge case}
What happens when you take $x=3,y=1,n=2$? The theorem breaks down actually.
This is because all odd squares are $1\pmod 8$ and thus
\[8\mid 3^2-1^2\]
while $\nu_2(3-1)$ is only $1$.
\end{frame}
\begin{frame}{Proof outline}
Should tackle the case when $n\neq 0\pmod p$ separately.
What happens when you take $v_p$ of
\[\frac{x^n-y^n}{x-y}=x^{n-1}+\dots+y^{n-1}?\]
We're supposed to get $0$- in other words, nonzero mod $p$.\\
However recall we assumed $p\mid x-y$ aka $y\equiv x \pmod p$.
Thus the thing on the right just becomes 
\[nx^{n-1}\]
which is obviously nonzero by definition.\\
So in other words, when not $p\mid n$ we have $\nu_p(x^n-y^n)=
\nu_p(x-y)$.
\end{frame}
\begin{frame}{General case}
We already tackled $p\not\mid n$ earlier so now we can break down $n$
as $p^{\nu_p n} n'$ for some integer $n'$ not a multiple of $p$.\\
\[\nu_p(x^n-y^n)=\nu_p(x^{n/n'}-y^{n/n'})
=\nu_p(x^{p^{\nu_p n}}-y^{p^{\nu_p n}}).\]
So in other words we only care about prime powers $n$: want 
\[\nu_p(x^{p^m}-y^{p^m})=m+\nu_p(x-y)\]
\end{frame}
\begin{frame}{Induction to finish}
It suffices to prove
\begin{block}{Claim}
$\nu_p(x^p-y^p)=1+\nu_p(x-y).$
\end{block}
\begin{proof}
Much of the same. We do have to take some care to obtain exactly $1$.
We let $y=x+kp$ and calculate mod $p^2$. In other words we want
\[x^{p-1}+\dots+ y^{p-1} \equiv smth\cdot p \pmod{p^2}.\]
We can just get the other terms by cutting the $p^{\ge2}$ terms 
automatically.\\
Expand $y^n=(x+kp)^n$ by binomial theorem:
\[x^n+nx^{n-1}(kp)+\dots\]
\end{proof}
\end{frame}
\begin{frame}{Finish inductive step}
Each of the terms is $(0\le a \le p-1)$ 
\[x^{n-1-a}y^a\equiv x^{n-1-a}(x^a+ax^{a-1}kp)\pmod{p^2}\]
\[Original \equiv \sum_{a=0}^{p-1}(x^{p-1}+ax^{p-2}kp)=\sum x^{p-1}
+kpx^{p-2}\sum a.\]
The latter thing is just zero so we end up with 
\[px^{p-1}\pmod{p^2}\]
which is a multiple of $p$ but not $p^2$ so we win.
\end{frame}
\begin{frame}{Solve AIME prob}
\begin{block}{Restate prob}
Let $n$ be the least positive integer for which 
$149^n - 2^n$ is divisible by $3^3 \cdot 5^5 \cdot 7^7$. 
Find the number of positive divisors of $n$.
\end{block}

\begin{itemize}
\item Need $\nu_3(149^n-2^n)\ge 3$. By LTE this is $\nu_3(149-2)+\nu_3 n$.
In other words we need $9$ to divide $n$.
\item Do the same thing for $7$. Then $2+\nu_7 n \ge 7$ so $7^5\mid n$.
\item When we do $p=5$, notice that not $5\mid 149-2$. It turns out that 
in order for 
\[5\mid 149^n-2^n\]
we need $4\mid n$. So if we let $n'=n/4$ then
\[5^5 \mid (149^4)^{n'}-(2^4)^{n'}\]
and by the exact same logic we get $5^4\mid n'$.
\end{itemize}
\end{frame}
\begin{frame}{Ans extract}
To summarize, all the above forces $3^2\cdot 7^5 \cdot 4\cdot 5^4$
to divide $n$ and that number is just the smallest $n$.
The answer ends up being $3\cdot 6 \cdot 3\cdot 5=\boxed{270}$.
\end{frame}
\end{document}
